\chapter{Discussion}
\label{sec:discussion}

This chapter interprets the experimental results in the context of the research
questions posed in Chapter~1, examines the failure modes of the proposed system,
compares performance with the most relevant prior work, and identifies the
limitations of the current evaluation.

\section{Classification Performance and Research Question}
\label{sec:disc_perf}

The central research question asked whether a context-aware, multi-cue heuristic
operating on geometry and facial landmarks could selectively anonymize bystanders
while preserving interaction partners, without biometric recognition or cloud
offloading. The Phase~2a evaluation results (PSR = 87.50\%, IPR = 84.62\%) confirm
that this is feasible using only bounding-box geometry. A PSR of 87.50\% means that
the system incorrectly exposes 12.5\% of bystander tracks; equivalently, 5 of the
40 bystander tracks in the evaluation dataset were mistakenly preserved unblurred.
An IPR of 84.62\% means that 15.38\% of interaction partners were unnecessarily
blurred; equivalently, 4 of 26 partner tracks were incorrectly anonymized.

Both error types carry concrete costs. A false positive (bystander preserved as
partner) constitutes a \emph{privacy leak}: the system fails to protect an
individual who did not consent to being recorded. A false negative (partner blurred
as bystander) constitutes a \emph{utility loss}: the wearer's conversation partner
is anonymized, defeating the purpose of selective anonymization. The observed ratio
of 5 FPs to 4 FNs reflects a roughly balanced trade-off, consistent with the
system's design objective of avoiding strong asymmetry in either direction.

The addition of Phase~2b landmark cues is expected to reduce false positives
specifically: the yaw veto correctly classifies faces that are turned significantly
away from the camera as bystanders, even if those faces happen to be large or
central in the frame (which would incorrectly pass the Phase~2a geometry rule). The
V-VAD boost conversely reduces false negatives by promoting faces that exhibit
sustained mouth movement to \emph{Interactor} status, even when centrality is
marginally insufficient.

The Phase~2b results (Table~\ref{tab:metrics_2b}) are consistent with the first
expectation: every bystander in the human-annotated sample was blurred, so no
bystander's face was left visible in the output (FP = 0, PSR = 100\%). This includes
eight bystanders who appeared in interaction clips next to an active partner. The
second expectation was not borne out: IPR was 57.14\%, with three of seven partners
blurred. As Section~\ref{sec:disc_failures} shows, these partners were not blurred
because the head-pose check judged them to be facing away. They were blurred because
their faces were too small in the frame to pass the proximity check, which is applied
before the head-pose and speaking checks are considered. In the
evaluated sample, the full system therefore behaved as a privacy-first classifier.

The larger video-clip evaluation confirms this pattern with much stronger evidence.
With almost perfect agreement between five annotators ($\kappa = 0.887$) and 102 clips,
the live Phase~2b system blurred 96.97\% of bystanders and preserved 56.52\% of
partners, and when it preserved a face it was correct 97.5\% of the time. The
privacy-first behaviour of the live system is therefore a stable property rather
than an artefact of a small sample. The offline second pass shows that most of the
remaining utility loss can be recovered without weakening privacy substantially:
by applying the latch retroactively and recognising the wearer's group, IPR rose to
78.26\% at a PSR of 93.94\%, giving the best balance of the configurations I
evaluated (F$_1$ = 86.40\%). Together, these results answer the research question
positively: interaction-aware selective anonymization can protect most bystanders
while preserving most of the people the wearer is with, without face recognition or
additional sensors.

\section{Comparison with BystandAR}
\label{sec:disc_bystandar}

The closest prior system, BystandAR \citep{bystandar}, achieved PSR = 98.14\% and
IPR = 96.27\% on a 16-participant evaluation with dedicated gaze-tracking and
directional microphone hardware. The proposed system achieves lower accuracy (PSR
87.50\%, IPR 84.62\% for Phase~2a; PSR 100\%, IPR 57.14\% for Phase~2b on a smaller
annotated sample) using only a monocular RGB camera and CPU-only inference. This
performance gap is expected and attributable to two factors: the richness of input
modalities and the evaluation scale.

BystandAR uses eye gaze as a direct signal of where the wearer is looking, which is
a near-perfect proxy for interaction intent. The proposed system infers this
indirectly from face geometry and pose, which is a weaker signal subject to ambiguity
when bystanders happen to be close to the camera wearer (for example, in a crowded
corridor). Additionally, BystandAR's evaluation involved 16 controlled participant
sessions with standardised interaction protocols, whereas the present evaluation used
21 naturalistic clips without controlled conditions.

Despite this gap, the proposed system's significance lies in its deployment
constraints: it requires no gaze tracker, no directional microphone array, and no
dedicated wearable hardware. It operates on a commodity RGB camera using only
CPU-bound inference, making it directly deployable on low-cost edge devices without
hardware modification. This positions it as a practical privacy layer for consumer
AR glasses such as the Meta Ray-Ban, which lack the dedicated sensors required by
BystandAR.

\section{Computational Feasibility and Optimisation Pathways}
\label{sec:disc_fps}

The measured throughput of 11.62 FPS on a laptop CPU shows that the complete
pipeline runs at interactive rates without any GPU acceleration. On the Raspberry Pi 5,
the projected throughput is approximately 3.3 FPS (Phase~2a) and 2.7 FPS (Phase~2b)
with the unoptimised PyTorch implementation. Since face detection accounts for 64\% of
the processing time, three optimisation pathways can raise this substantially:

\begin{enumerate}
  \item \textbf{ONNX export and quantisation.} Exporting YOLOv8n-face to ONNX INT8
        format and running inference via ONNXRuntime with NEON-optimised kernels for
        ARM can reduce detection latency by a factor of 2-3$\times$ on ARM hardware
        relative to the PyTorch runtime used in the current evaluation.

  \item \textbf{Reduced input resolution.} Processing frames at $416 \times 416$
        rather than $640 \times 640$ reduces YOLO inference time by approximately
        35\% with a modest reduction in detection recall for small, distant faces.
        For the target use case of close-range egocentric video, this trade-off is
        likely acceptable.

  \item \textbf{Landmark extraction batching.} In the current implementation,
        MediaPipe Face Landmarker is called once per face per frame. Batching
        multiple face crops into a single inference call would reduce per-face
        overhead, particularly when multiple faces are present simultaneously.
\end{enumerate}

With all three optimisations applied, achieving 8-12~FPS on the Raspberry Pi 5 is
plausible, which is sufficient for processing recorded egocentric video (a common
use case in AR privacy logging). Full-frame-rate operation on the Pi~5 can be reached by
running the detector on a neural processing unit (NPU) or AI co-processor, which
is included in many current smart-glasses processors.

\paragraph{Comparison with commercial smart-glasses hardware.}
To judge whether the evaluation hardware is representative, I compared it with
current camera-equipped glasses (Table~\ref{tab:glasses}). Ray-Ban Meta and Meta
Ray-Ban Display glasses use the Qualcomm Snapdragon AR1 Gen~1 platform
\citep{ref:qcom_ar1}, which combines a quad-core CPU (reported at up to about
1.9~GHz), dual image signal processors and a dedicated Hexagon neural processing unit
(NPU), with 2~GB of RAM in the Display model \citep{ref:gsmarena_display}. Snap
Spectacles (fifth generation) split the work across two Qualcomm Snapdragon processors
and run for about 45 minutes per charge \citep{ref:spectacles5}, and Meta's Aria
Gen~2 research glasses perform their machine perception on a custom low-power
co-processor \citep{ref:aria2}.

\begin{table}[htbp]
\begin{singlespace}
\centering
\caption{Evaluation hardware compared with current camera-equipped glasses.
Values are from manufacturer and press sources; ``not published'' means the
manufacturer has not released the figure.}
\label{tab:glasses}
\renewcommand{\arraystretch}{1.3}
\small
\setlength{\tabcolsep}{4pt}
\begin{tabular}{@{}>{\raggedright\arraybackslash}p{0.20\linewidth}
                 >{\raggedright\arraybackslash}p{0.21\linewidth}
                 >{\raggedright\arraybackslash}p{0.16\linewidth}
                 >{\raggedright\arraybackslash}p{0.15\linewidth}
                 >{\raggedright\arraybackslash}p{0.20\linewidth}@{}}
\toprule
\textbf{Device} & \textbf{Processor} & \textbf{AI accelerator} & \textbf{Battery} & \textbf{Where vision AI runs} \\
\midrule
Laptop (my evaluation)    & Intel Core i5-13420H & none used & mains / battery & on device (CPU) \\
Raspberry Pi 5 (target) & 4$\times$ Cortex-A76, 2.4~GHz & none & external & on device (CPU) \\
Ray-Ban Meta          & Snapdragon AR1 Gen~1 & Hexagon NPU & about 4~h & mostly phone and cloud \\
Meta Ray-Ban Display  & Snapdragon AR1 Gen~1 & Hexagon NPU & about 6~h & mostly phone and cloud \\
Snap Spectacles (5th gen.) & 2$\times$ Snapdragon & not published & about 45~min & on device \\
Meta Aria Gen~2       & custom co-processor & custom & 6-8~h & on device \\
\bottomrule
\end{tabular}
\end{singlespace}
\end{table}

This comparison supports the design choices of my research in three ways. First, the
Raspberry Pi 5 is a reasonable stand-in for the CPU of current glasses: its four
Cortex-A76 cores at 2.4~GHz are of the same class as the AR1's quad-core CPU, so
the projected Pi~5 figures are a realistic estimate of CPU-only performance on such
devices. Second, my evaluation used the CPU only, whereas every current glasses
platform includes a neural processing unit. Face detection, which accounts for 64\%
of the pipeline's processing time, is exactly the kind of workload an NPU
accelerates. On a Snapdragon smartphone NPU, a quantised YOLO detector runs in about
5.6~ms per frame against 53~ms on the same phone's CPU \citep{ref:ultralytics_qnn},
and Qualcomm's quantised face detector runs in under 1~ms \citep{ref:qcom_facedet}.
Glasses processors are smaller than smartphone processors, so these figures are not
directly transferable, but they indicate that the CPU-only results reported here are
a worst case for real glasses. Third, battery capacity, rather than raw speed, is
the main constraint on glasses: Spectacles run for about 45 minutes, and current
Ray-Ban Meta glasses send most vision AI to the paired phone or the cloud. Because
my pipeline keeps all face processing on the device, it avoids sending bystanders'
faces off the glasses, which Meta itself identifies as a central privacy concern for
camera glasses \citep{ref:meta_bystander}. Reducing how often the detector runs, as
proposed above, also directly reduces energy use.

\section{Failure Mode Analysis}
\label{sec:disc_failures}

Analysis of the five false-positive tracks (bystanders incorrectly preserved as
partners) reveals a consistent pattern: all five occurred in the \emph{interaction}
video category, where people standing near the camera wearer were never directly
addressed but occupied a large, central position in the frame for an extended
period. In these cases, the area ratio and centrality thresholds are satisfied
because the individual is physically close to the wearer, even though no interaction
occurs. This ambiguity is precisely the scenario that the Phase~2b yaw veto is
designed to address: a person who is nearby but looking in a different direction
(yaw $> 40^\circ$) should be vetoed as a bystander regardless of proximity.

The four false-negative tracks (partners incorrectly blurred) arose from brief
interactions that began after a long period in which the face was peripheral, so
the temporal hysteresis had not yet committed the track to the \emph{Interactor}
state.

\paragraph{Phase 2b failure modes.}
All three Phase~2b false negatives share a single cause: at the annotated frame,
each partner's face occupied between 0.36\% and 0.72\% of the frame area, below the
area threshold $\tau_{\mathrm{area}} = 1.2\%$. Two of them occur in one frame of a
seated group setting, where the same person was covered by two overlapping
detections that SORT tracked as separate tracks; this duplicate-track artefact
counts one missed partner twice. The third is a partner standing at a service
counter near the edge of the frame, small and partly turned (yaw $\approx 25^\circ$,
below the $40^\circ$ veto). In each case the area gate rejected the face before any
landmark cue was considered. Because the V-VAD boost also requires
$r_{\mathrm{area}} \geq \tau_{\mathrm{area}}$, it cannot rescue partners who converse
from further away than the area threshold assumes, for example across a table.

A further factor limits the influence of the landmark cues. MediaPipe Face
Landmarker returned landmarks for only 50.8\% of the 224{,}519 logged face
detections. For the remainder, which are typically small, motion-blurred or partially
occluded faces, no yaw or MAR value is available and the classifier falls back to the
Phase~2a geometry rule. Both observations point to the same improvement: relaxing
the area threshold when landmarks confirm a frontal, speaking face, and improving
landmark coverage on small faces (for example by upscaling the crop before landmark
extraction).

\paragraph{Video-clip failure modes.}
The video-clip evaluation confirms the same dominant cause on a larger sample. Of
the 30 partners the live system blurred, 25 had a median face size below the area
threshold. The median face size of missed partners was 0.91\% of the frame, against
1.72\% for partners the system preserved. Most of these were friends sitting further
from the wearer around a table, who are partners under the group definition but
appear smaller than the threshold assumes. In addition, 10 of the 30 were confirmed
as partners later in the same video, which shows that the live system often
recognises a partner correctly but too late. The offline second pass addresses both
causes directly, and recovered 15 of the 30. The only bystander the live system
preserved was left unblurred for exactly half of the clip, the threshold of the
per-clip decision rule.

The video-clip study also gives a direct view of detector errors. Of the 104
sampled tracks, two (1.9\%) were not people: a face printed on a magazine cover,
which the system blurred, and a product package on a shop shelf, which the system
classified as a partner and left unblurred. Neither error affects privacy, since no
real person is involved, but the second shows that the interaction cues are applied
to any detection that persists in the centre of the view. This detection had a
median confidence of only 0.53, just above the detection threshold of 0.45, and the
landmark extractor found a face in only 1 of its 73 frames. Requiring a higher
detection confidence, or successful landmark extraction, before a track can be
confirmed as a partner would filter out such non-face detections.

\section{Limitations}
\label{sec:disc_limitations}

The following limitations should be considered when interpreting the results of this
evaluation.

\textbf{Dataset scale and diversity.} The evaluation dataset comprises 21 clips
collected from online sources rather than recorded for this study, so capture
conditions such as camera model, mounting position and frame rate are uncontrolled
and vary between clips. Interaction was evaluated only indoors, and the interaction
clips are dominated by seated tabletop scenes; partner detection was therefore not
tested in outdoor conversations, while bystander behaviour was tested only outdoors.
Results may not generalise to other settings, lighting conditions or interaction
styles. A larger dataset in which both categories span the same range of
environments would provide a more robust assessment of generalisation.

\textbf{Annotation scale and protocol.} I assigned the Phase~2a ground truth
alone. Phase~2b was evaluated twice with blind protocols: on 29 faces from still
frames by four annotators (mean $\kappa = 0.587$), and on 102 video clips by five
annotators (Fleiss' $\kappa = 0.887$). I was one of the annotators in both studies
and was familiar with the system design, although my labels agreed with the
independent annotators' labels in 92-98\% of the video clips. The two Phase~2b
studies also used different partner definitions (actively talking to the wearer,
and member of the wearer's group). Because Phase~2a and Phase~2b were evaluated on
different samples, differences between them are indicative rather than conclusive.

\textbf{Offline second pass.} The retroactive latch and the friend-group rule were
designed after examining the errors of the video-clip evaluation and were evaluated
on the same labels. Although the result was stable across a wide range of time
limits, a fresh set of labelled clips is needed for an independent test. Both rules
also need the whole recording, so they apply to recorded video only.

\textbf{Heuristic thresholds.} I set all eight threshold parameters
(Table~\ref{tab:params}) empirically on the evaluation dataset itself. This risks
overfitting to the specific conditions of the collected clips. A held-out validation set or
cross-validation procedure would provide a fairer estimate of threshold robustness.

\textbf{Absence of audio.} The V-VAD cue relies entirely on lip movement as a
speech proxy and cannot distinguish between a person mouthing words, laughing,
chewing, or speaking. Incorporating even a rudimentary microphone-based voice
activity detector as a secondary cue would substantially reduce V-VAD ambiguity at
negligible hardware cost on most wearable devices.

\textbf{Hardware projection.} The Raspberry Pi 5 throughput figures (approximately
3.3 FPS for Phase~2a and 2.7 FPS for Phase~2b) are estimates derived from a scaling factor rather than a direct
measurement. They should be treated as an order-of-magnitude projection; actual
performance may differ depending on OS configuration, thermal throttling, and the
version of underlying inference libraries installed. The battery-power measurements
are single runs affected by rising processor temperature, and I did not measure
energy consumption directly, so the battery cost of continuous operation on glasses
remains to be quantified.
